\setcounter{section}{1}

  \zwischenueberschrift{Soal latihan}

\inputaufgabegibtloesung
{}
{

Misalkan

\mavergleichskette
{\vergleichskette
{ G
}
{ \subseteq }{ \R^n
}
{  }{
}
{    }{
}
{   }{
}
}
{}{}{}
\definitionsverweis {terbuka}{}{}
dan
\maabbdisp {\varphi} { G } { \R^m
} {}
sebuah
\definitionsverweis {pemetaan terdiferensialkan secara kontinu}{}{,}
yang di titik

\mavergleichskette
{\vergleichskette
{ P
}
{ \in }{ G
}
{  }{
}
{    }{
}
{   }{
}
}
{}{}{}
mempunyai
\definitionsverweis {diferensial total}{}{}
yang surjektif. Misalkan

\mavergleichskette
{\vergleichskette
{ v
}
{ \in }{ T_PZ
}
{  }{
}
{    }{
}
{   }{
}
}
{}{}{}
sebuah vektor di
\definitionsverweis {ruang tangen pada serat}{}{}
$Z$ dari $\varphi$ yang melalui $P$. Buktikan bahwa terdapat
\definitionsverweis {kurva terdiferensialkan secara kontinu}{}{}
\maabbdisp {\gamma} { ]- \epsilon, \epsilon[ } { Z
} {}
\zusatzklammer {untuk suatu

\mavergleichskettek
{\vergleichskettek
{ \epsilon
}
{ > }{ 0
}
{  }{
}
{    }{
}
{   }{
}
}
{}{}{}
yang sesuai} {} {}
dengan

\mavergleichskette
{\vergleichskette
{ \gamma(0)
}
{ = }{ P
}
{  }{
}
{    }{
}
{   }{
}
}
{}{}{}
dan

\mavergleichskettedisp
{\vergleichskette
{ \gamma'(0)
}
{ =} { v
}
{ } {
}
{ } {
}
{ } {
}
}
{}{}{.}

}
{} {}

\inputaufgabe
{}
{

Misalkan

\mavergleichskette
{\vergleichskette
{ W
}
{ \subseteq }{ \R^n
}
{  }{
}
{    }{
}
{   }{
}
}
{}{}{}
\definitionsverweis {terbuka}{}{,}
\maabb {h} { W } { \R
} {}
sebuah
\definitionsverweis {fungsi terdiferensialkan secara kontinu}{}{}
dan

\mavergleichskette
{\vergleichskette
{ Y
}
{ = }{ h^{-1}(0)
}
{  }{
}
{    }{
}
{   }{
}
}
{}{}{}
merupakan
\definitionsverweis {serat}{}{}
di atas $0$, dengan $h$
\definitionsverweis {reguler}{}{}
di setiap titik pada $Y$. Misalkan
\maabbdisp {g} { W } {\R
} {}
sebuah fungsi terdiferensialkan secara kontinu yang tidak mempunyai titik nol pada $Y$. Buktikan bahwa $Y$ juga dapat diperoleh sebagai serat dari

\mavergleichskettedisp
{\vergleichskette
{ \tilde{h}
}
{ =} { g h
}
{ } {
}
{ } {
}
{ } {
}
}
{}{}{,}
bahwa
\definitionsverweis {gradien}{}{}
dari $h$ dan $\tilde{h}$ merupakan kelipatan skalar satu sama lain, dan bahwa
\definitionsverweis {ruang-ruang tangen}{}{}
hanya bergantung pada $Y$.

}
{} {}

\inputaufgabe
{}
{

Misalkan

\mavergleichskette
{\vergleichskette
{ W
}
{ \subseteq }{ \R^n
}
{  }{
}
{    }{
}
{   }{
}
}
{}{}{}
\definitionsverweis {terbuka}{}{,}
\maabb {h} { W } { \R
} {}
sebuah
\definitionsverweis {fungsi terdiferensialkan secara kontinu}{}{}
dan

\mavergleichskette
{\vergleichskette
{ Y
}
{ = }{ h^{-1}(c)
}
{  }{
}
{    }{
}
{   }{
}
}
{}{}{}
merupakan
\definitionsverweis {serat}{}{}
di atas

\mavergleichskette
{\vergleichskette
{ c
}
{ \in }{ \R
}
{  }{
}
{    }{
}
{   }{
}
}
{}{}{,}
dengan $h$
\definitionsverweis {reguler}{}{}
di setiap titik pada $Y$. Misalkan
\maabbdisp {L} {\R^n } { \R^n
} {}
sebuah
\definitionsverweis {isomorfisme linear}{}{,}

\mavergleichskette
{\vergleichskette
{ W'
}
{ = }{ L^{-1}(W)
}
{  }{
}
{    }{
}
{   }{
}
}
{}{}{}
dan

\mavergleichskettedisp
{\vergleichskette
{ Z
}
{ =} { L^{-1}(Y)
}
{ =} { (h \circ L)^{-1}(c)
}
{ } {
}
{ } {
}
}
{}{}{.}
Buktikan bahwa konsep
\definitionsverweis {ruang tangen}{}{,}
medan vektor tangensial, dan solusi persamaan diferensial tangensial yang terkait pada $Y$ dan $Z$ secara umum saling bersesuaian
\zusatzklammer {yakni dapat dipetakan satu sama lain dengan bantuan $L$} {} {.}
Bagaimana dengan gradiennya?

}
{} {}

\inputaufgabe
{}
{

Misalkan

\mavergleichskette
{\vergleichskette
{ Y
}
{ = }{ h^{-1}(c)
}
{  }{
}
{    }{
}
{   }{
}
}
{}{}{}
sebuah
\definitionsverweis {hipermuka terdiferensialkan}{}{}
dan misalkan
\maabbdisp {f} { U} { \R
} {}
sebuah fungsi terdiferensialkan secara kontinu yang didefinisikan pada suatu lingkungan terbuka

\mavergleichskette
{\vergleichskette
{ Y
}
{ \subseteq }{ U
}
{  }{
}
{    }{
}
{   }{
}
}
{}{}{}
Andaikan $f$ di

\mavergleichskette
{\vergleichskette
{ P
}
{ \in }{ Y
}
{  }{
}
{    }{
}
{   }{
}
}
{}{}{}
mempunyai
\definitionsverweis {ekstremum lokal}{}{}
pada $Y$. Buktikan bahwa
\definitionsverweis {komponen tangensial}{}{}
dari $\operatorname{Grad} \,  f  (  P )$ sama dengan $0$.

}
{} {}

\inputaufgabe
{}
{

Misalkan

\mavergleichskette
{\vergleichskette
{ Y
}
{ \subseteq }{ \R^3
}
{  }{
}
{    }{
}
{   }{
}
}
{}{}{}
adalah permukaan bola satuan. Tentukan
\definitionsverweis {dekomposisi ortogonal}{}{}
vektor

\mavergleichskette
{\vergleichskette
{ \begin{pmatrix}  6  \\ 7 \\  -3   \end{pmatrix}
}
{ \in }{ \R^3
}
{  }{
}
{    }{
}
{   }{
}
}
{}{}{}
menjadi komponen tangensial dan komponen ortogonal di titik

\mavergleichskette
{\vergleichskette
{ \begin{pmatrix}  -  { \frac{   1  }{  2 } }  \\ { \frac{   1  }{  2 } } \\  { \frac{   1  }{  \sqrt{2}  } }   \end{pmatrix}
}
{ \in }{ Y
}
{  }{
}
{    }{
}
{   }{
}
}
{}{}{.}

}
{} {}

\inputaufgabe
{}
{

Misalkan

\mavergleichskette
{\vergleichskette
{ Y
}
{ \subseteq }{ \R^3
}
{  }{
}
{    }{
}
{   }{
}
}
{}{}{}
adalah
\definitionsverweis {permukaan bola satuan}{}{.}
Misalkan
\maabbdisp {\gamma} {[0, 2 \pi]} { Y
} {}
merupakan gerak melingkar beraturan yang mengelilingi permukaan bola pada ketinggian ${ \frac{   1  }{  2 } }$.
\aufzaehlungzwei {Parametrisasikan $\gamma$.
} {Tentukan
\definitionsverweis {percepatan tangensial}{}{}
dari $\gamma$ pada setiap saat

\mavergleichskette
{\vergleichskette
{ t
}
{ \in }{ [0, 2 \pi]
}
{  }{
}
{    }{
}
{   }{
}
}
{}{}{.}
}

}
{} {}

\inputaufgabe
{}
{

Misalkan $Y$ adalah
\definitionsverweis {grafik}{}{}
dari
\maabbeledisp {f} {\R^2} { \R
} {(x,y)} { x^2+y^3
} {,}
dan misalkan
\maabbdisp {\gamma} {\R} { Y
} {}
adalah lintasan pada grafik tersebut yang berada di atas kurva dasar
\mathl{t \mapsto (t,t)}{.} Tentukan, untuk setiap

\mavergleichskette
{\vergleichskette
{ t
}
{ \in }{ \R
}
{  }{
}
{    }{
}
{   }{
}
}
{}{}{,}
\definitionsverweis {percepatan tangensial}{}{}
dari $\gamma$ di $\gamma(t)$.

}
{} {}

\inputaufgabe
{}
{

Misalkan

\mavergleichskette
{\vergleichskette
{ Y
}
{ \subseteq }{ \R^n
}
{  }{
}
{    }{
}
{   }{
}
}
{}{}{}
sebuah hiperbidang, yaitu serat dari fungsi linear

\mavergleichskettedisp
{\vergleichskette
{ h(x_1 , \ldots , x_n)
}
{ =} { a_1x_1 + \cdots + a_nx_n
}
{ } {
}
{ } {
}
{ } {
}
}
{}{}{}
di atas

\mavergleichskette
{\vergleichskette
{ c
}
{ \in }{ \R
}
{  }{
}
{    }{
}
{   }{
}
}
{}{}{}
dengan

\mavergleichskette
{\vergleichskette
{ a_i
}
{ \neq }{ 0
}
{  }{
}
{    }{
}
{   }{
}
}
{}{}{}
untuk setidaknya satu indeks $i$. Tafsirkan pengertian
\definitionsverweis {diferensial total}{}{,}
\definitionsverweis {gradien}{}{,}
\definitionsverweis {ruang tangen}{}{,}
dekomposisi menjadi komponen tangensial dan normal,
\definitionsverweis {percepatan tangensial}{}{,}
serta
\definitionsverweis {kurva geodesik}{}{}
dalam kasus khusus ini.

}
{} {}

\inputaufgabe
{}
{

Misalkan

\mavergleichskette
{\vergleichskette
{ Y
}
{ = }{ h^{-1}(c)
}
{  }{
}
{    }{
}
{   }{
}
}
{}{}{}
sebuah
\definitionsverweis {hipermuka terdiferensialkan}{}{,}
yang memuat garis
\mathl{P+ \R v}{.} Buktikan bahwa garis yang dilalui secara seragam tersebut merupakan
\definitionsverweis {kurva geodesik}{}{}
pada $Y$.

}
{} {}

Dua titik

\mathbed {P,Q \in K} {}
 {P \neq Q} {}
 {} {} {} {,}
disebut \stichwort {antipodal} {,} jika garis yang menghubungkannya melalui pusat bola.

\inputaufgabe
{}
{

Misalkan

\mavergleichskette
{\vergleichskette
{ K
}
{ \subseteq }{ \R^3
}
{  }{
}
{    }{
}
{   }{
}
}
{}{}{}
sebuah permukaan bola. Buktikan bahwa setiap dua
\definitionsverweis {titik yang tidak antipodal}{}{}

\mathbed {P,Q \in K} {}
 {P \neq Q} {}
 {} {} {} {,}
terletak pada tepat satu
\definitionsverweis {lingkaran besar}{}{}
dari $K$.

}
{} {}

\bild{ \begin{center}
\includegraphics[width=5.5cm]{ \bildeinlesungsvg {Planned flight map of the Oiseau Blanc} {svg} }
\end{center}
\bildtext {Mengapa pesawat terbang menempuh lintasan melengkung?} }

\bildlizenz { Planned flight map of the Oiseau Blanc.svg } {} {Pethrus} {Commons} {CC BY-SA 3.0} {}

\inputaufgabe
{}
{

Sebuah pesawat akan terbang dari Osnabrück menuju suatu tempat tujuan di belahan Bumi selatan. Mungkinkah rutenya lebih pendek jika pesawat itu mula-mula terbang ke arah utara?

}
{} {}

\inputaufgabe
{}
{

Lucy Sonnenschein merasa bahwa hari-hari terlalu singkat. Karena itu, setiap hari dan selama masih terang, ia memutuskan terbang ke arah barat melintasi satu zona waktu agar setiap harinya berlangsung $25$, bukan $24$, jam.
\aufzaehlungfuenf{Dapatkah Lucy Sonnenschein meningkatkan proporsi jam terang dalam hidupnya dengan strategi ini?
}{Dapatkah Lucy Sonnenschein memperpanjang hidupnya dengan strategi ini?
}{Sekarang Lucy mempunyai teman di seluruh dunia. Setelah berapa hari ia bertemu dengan mereka lagi?
}{Apa yang menyebabkan Lucy mengalami kehilangan waktu yang mengimbangi tambahan waktu hariannya?
}{Apakah hal ini berkaitan dengan teori relativitas?
}

}
{} {}

\inputaufgabe
{}
{

Berikan contoh kurva dua kali terdiferensialkan yang bergerak pada suatu
\definitionsverweis {lingkaran besar}{}{}
dari permukaan bola satuan, tetapi bukan
\definitionsverweis {kurva geodesik}{}{.}

}
{} {}

\inputaufgabe
{}
{

Sebuah \stichwort {segitiga geodesik} {} pada permukaan bola

\mavergleichskettedisp
{\vergleichskette
{ S^2
}
{ =} { \left\{ (x,y,z) \in \R^3  \mid   x^2+y^2+z^2 = 1        \right\}
}
{ } {
}
{ } {
}
{ } {
}
}
{}{}{}
terdiri atas tiga titik berbeda yang tidak terletak pada satu lingkaran besar,

\mavergleichskette
{\vergleichskette
{ A,B,C
}
{ \in }{ S^2
}
{  }{
}
{    }{
}
{   }{
}
}
{}{}{,}
dengan setiap pasangan titik dihubungkan oleh busur yang lebih pendek dari suatu
\definitionsverweis {lingkaran besar}{}{}
\zusatzklammer {setiap pasangan titik tidak antipodal} {} {.}
Buktikan bahwa jumlah sudut dalam segitiga itu lebih besar daripada $\pi$.

}
{} {}

  \zwischenueberschrift{Soal untuk dikumpulkan}

\inputaufgabe
{3}
{

Misalkan $Y$ adalah hipermuka yang diberikan oleh syarat

\mavergleichskettedisp
{\vergleichskette
{ 2x^2 +3y^2+5z^2
}
{ =} { 10
}
{ } {
}
{ } {
}
{ } {
}
}
{}{}{}
di $\R^3$. Tentukan dekomposisi ortogonal vektor

\mavergleichskette
{\vergleichskette
{ \begin{pmatrix}  4  \\ 2 \\  -5   \end{pmatrix}
}
{ \in }{ \R^3
}
{  }{
}
{    }{
}
{   }{
}
}
{}{}{}
menjadi komponen tangensial dan komponen ortogonal di titik

\mavergleichskette
{\vergleichskette
{ \begin{pmatrix}  -1 \\ -1\\  1   \end{pmatrix}
}
{ \in }{ Y
}
{  }{
}
{    }{
}
{   }{
}
}
{}{}{.}

}
{} {}

\inputaufgabe
{4}
{

Misalkan $Y$ adalah
\definitionsverweis {grafik}{}{}
dari
\maabbeledisp {f} {\R^2} { \R
} {(x,y)} { x^2+y^2
} {,}
dan misalkan
\maabbdisp {\gamma} {\R} { Y
} {}
adalah lintasan pada grafik tersebut yang berada di atas kurva dasar
\mathl{x \mapsto (x,1)}{.} Tentukan, untuk setiap

\mavergleichskette
{\vergleichskette
{ x
}
{ \in }{ \R
}
{  }{
}
{    }{
}
{   }{
}
}
{}{}{,}
\definitionsverweis {percepatan tangensial}{}{}
dari $\gamma$ di $\gamma(x)$.

}
{} {}

\inputaufgabe
{4}
{

Pada 11 April 2023 pukul 17.45, secara tiba-tiba dan mengejutkan semua orang, rotasi Bumi pada porosnya dihentikan. Pada saat yang sama, hambatan udara dan semua kerugian akibat gesekan ditiadakan. Semua manusia, hewan, dan benda yang tidak terpasang tetap atau tidak sempat berpegangan akan meneruskan gerak sesaatnya sesuai hukum inersia. Gravitasi tetap bekerja sehingga, untungnya, mereka tidak terlepas ke angkasa. Kapan penghapus papan tulis dari ruang kuliah di Osnabrück
\zusatzklammer {penghapus itu terbang keluar melalui jendela} {} {}
untuk pertama kalinya kembali ke posisi awalnya?

}
{} {}

\inputaufgabe
{4 (3+1)}
{

Misalkan

\mavergleichskette
{\vergleichskette
{ W
}
{ \subseteq }{ \R^n
}
{  }{
}
{    }{
}
{   }{
}
}
{}{}{}
\definitionsverweis {terbuka}{}{,}
\maabb {h} { W } { \R
} {}
sebuah
\definitionsverweis {fungsi terdiferensialkan secara kontinu}{}{}
dan

\mavergleichskette
{\vergleichskette
{ Y
}
{ = }{ h^{-1}(c)
}
{  }{
}
{    }{
}
{   }{
}
}
{}{}{}
merupakan
\definitionsverweis {serat}{}{}
di atas

\mavergleichskette
{\vergleichskette
{ c
}
{ \in }{ \R
}
{  }{
}
{    }{
}
{   }{
}
}
{}{}{,}
dengan $h$
\definitionsverweis {reguler}{}{}
di setiap titik pada $Y$. Misalkan
\maabbdisp {\gamma} { I } { Y
} {}
sebuah kurva dua kali
\definitionsverweis {terdiferensialkan secara kontinu}{}{.}
\aufzaehlungzwei {Buktikan bahwa jika $\gamma$ merupakan
\definitionsverweis {kurva geodesik}{}{,}
maka
\definitionsverweis {norma}{}{}
dari $\gamma'$ adalah konstan.
} {Berikan contoh kurva dengan norma kecepatan konstan yang bukan kurva geodesik.
}

}
{} {}

\inputaufgabe
{3}
{

Misalkan
\maabb {h,g} {\R^n } { \R
} {}
adalah
\definitionsverweis {fungsi-fungsi terdiferensialkan secara kontinu}{}{}
dan

\mavergleichskette
{\vergleichskette
{ Y
}
{ = }{ h^{-1}(c)
}
{  }{
}
{    }{
}
{   }{
}
}
{}{}{}
serta

\mavergleichskette
{\vergleichskette
{ Z
}
{ = }{ g^{-1}(d)
}
{  }{
}
{    }{
}
{   }{
}
}
{}{}{}
adalah
\definitionsverweis {hipermuka terdiferensialkan}{}{}
yang terkait, dengan $h$ di setiap titik pada $Y$ dan $g$ di setiap titik pada $Z$ bersifat
\definitionsverweis {reguler}{}{.}
Misalkan

\mavergleichskette
{\vergleichskette
{ T
}
{ = }{ Y \cap Z
}
{  }{
}
{    }{
}
{   }{
}
}
{}{}{}
dan misalkan
\maabbdisp {\gamma} { I } { T
} {}
sebuah kurva dua kali terdiferensialkan secara kontinu yang, jika dipandang sebagai kurva pada $Y$, merupakan
\definitionsverweis {kurva geodesik}{}{.}
Apakah $\gamma$ juga merupakan kurva geodesik pada $Z$?

}
{} {}
